% =====================================================================
% CHAPTER 3: Methodology  (~22 pages)
% WRITING SESSION INSTRUCTIONS (keep until the chapter is done):
%   - Read thesis/WRITING_PLAN.md section 5 FIRST.
%   - Single source of truth: model/designs/StereoLite_yolo_ctx_gev4/
%     FINAL_MODEL_ARCHITECTURE.md (sections 3-11 blocks, 12 loss, 16
%     metrics, 18 pre-vetted claim wording). Cross-check against the
%     TRAINED code in model/designs/StereoLite_yolo_ctx_gev4_opt/.
%   - Register: method sections ASSERT in flat present tense; tensor
%     dims inline in prose; every equation gets a purpose clause before
%     and a where-clause after (patterns.md section 6).
%   - The equations below are PRE-PLACED and verified against the arch
%     doc; do not renumber or hand-number them; reference with \eqref.
%   - Model name: "StereoLite" thesis-wide; the internal variant slug
%     (gev4_opt_narrow_plane) appears only in Appendix A.
%   - POs: PO2 (formulation), PO3 (design, 3.2-3.9 + 3.12), PO5 (modern
%     tools, 3.14). CEP K5/K8 reinforced.
% =====================================================================
\chapter{Methodology}

\section{Problem Formulation and Design Requirements}
% Sources: WRITING_PLAN 5.1 row 3.1; THESIS_PLAN envelope memory.
% Content: rectified pair -> dense disparity; design requirements as the
% hard envelope (< 3 M params, real-time on 4 to 6 TOPS, < 200 MB
% inference memory, INT8-exportable ops); research questions RQ1
% (accuracy at budget), RQ2 (zero-shot generalization), RQ3 (edge
% latency), RQ4 (tolerance to imperfect rectification, objective 2);
% assumptions (rectified input, max disparity 192 px).
Stereo matching takes a pair of rectified images and produces a dense
disparity map, from which metric depth follows directly through the
calibrated relation of \eqref{eq:depth_from_disparity}. This thesis
addresses that task under a hardware budget rather than an accuracy
budget alone. The design requirements follow from the class of edge
accelerators found on mobile robots and drones: the network must be
compact enough for on-device memory, must run in real time on a device
of roughly four to six tera-operations per second, must keep a small
inference working set, and must use operators that quantize cleanly to
eight-bit integer arithmetic. These requirements frame four research
questions. The first asks whether a network within a few million
parameters can reach competitive in-domain accuracy on synthetic data.
The second asks whether it generalizes zero-shot to an unseen real
domain. The third asks whether it meets the real-time latency target on
the intended device class. The fourth asks whether it tolerates the
small vertical misalignments that arise when rectification is imperfect,
which is the second objective of this thesis. We assume rectified
binocular input and a maximum disparity of 192 pixels at the inference
resolution, and we train on synthetic data with dense ground truth.

\section{System Overview}
% Narrate the full pipeline once against Fig 3.1; one paragraph per
% stage; name the inherited-and-modified mechanisms with their Ch2
% ancestry: group-wise cost (\eqref{eq:gwc_lit}), ConvGRU + convex
% upsampling (RAFT-Stereo), GEV (IGEV), tile-plane hypotheses (HITNet).
Figure \ref{fig:architecture} gives the complete pipeline. The proposed
model, StereoLite, is a coarse-to-fine tile-hypothesis network with light
recurrent refinement. It inherits three mechanisms from the iterative and
3D-cost-volume literature and adapts each to the edge budget. The matching
cost is a group-wise correlation volume in the manner of GwcNet
\cite{guo2019gwcnet}, restricted to a single coarse scale. The refinement
operator is a convolutional gated recurrent unit and the final upsampling
is a learned convex combination, both adapted from RAFT-Stereo
\cite{lipson2021raftstereo} but run for only a handful of iterations at
coarse resolution. A quarter-resolution geometry encoding volume, in the
spirit of IGEV-Stereo \cite{xu2023igev}, supplies a global disparity
proposal that a learned gate blends into the recurrent estimate. Slanted
tile planes, following the HITNet paradigm \cite{tankovich2021hitnet},
carry a local surface hypothesis across scales so that disparity remains
sub-pixel without rebuilding a cost volume at every level. A shared
feature encoder feeds the matching path, and a separate left-image
context encoder feeds the recurrent state.

\begin{figure}[!htbp]
\centering
% INTERIM: viewer render of figures/fig_3_1_architecture.drawio. Once
% hand-adjustments in draw.io are done, export a cropped PDF over
% fig_3_1_architecture.pdf and swap the filename below.
\includegraphics[width=\textwidth]{fig_3_1_architecture_preview}
\caption{\textbf{Overview of the proposed StereoLite architecture.} A shared feature encoder (1) feeds a group-wise cost volume whose 3-D regularization initializes the slanted tile state (2); convolutional GRU updates refine the state across three scales (3) with plane-aware propagation between them; a narrow-band geometry encoding volume corrects the estimate through a learned fail-soft gate (4); plane rendering with learned convex upsampling recovers the full-resolution disparity (5). Red dots mark supervised outputs.}
\label{fig:architecture}
\end{figure}

Table \ref{tab:scales} summarizes the spatial sizes and the main
operation at each scale of the pipeline for the default 384 by 640
input.

\begin{table}[!htbp]
\centering
\caption{\textbf{Scales of the pipeline.} Spatial sizes at the default $384\times 640$ input, the encoder channel count at each scale, and the main operation performed there.}
\label{tab:scales}
\small
\begin{tabular}{lccl}
\toprule
Scale & Spatial size & Channels & Main operation\\
\midrule
Full & $384\times 640$ & 3 & input and output\\
1/2 & $192\times 320$ & 32 & convex-upsample guidance\\
1/4 & $96\times 160$ & 128 & GEV fusion and 3 refinements\\
1/8 & $48\times 80$ & 256 & 3 refinements\\
1/16 & $24\times 40$ & 256 & cost initialization and 2 refinements\\
\bottomrule
\end{tabular}
\end{table}

\section{Shared Feature Encoder and Context Encoder}
% Source: FINAL_MODEL_ARCHITECTURE sections 4 + 5. YOLO26s layers 0-6,
% shared weights via batch concatenation; f4 128 / f8 256 / f16 256 ch;
% the separate left-only context encoder (GhostConv + SE, 32 ch out) and
% WHY matching features must not double as recurrent context.
The matching network is a truncated YOLO26s backbone whose first seven
blocks are retained, producing features at one quarter, one eighth, and
one sixteenth resolution with 128, 256, and 256 channels. A detection
backbone is a deliberate choice rather than a convenience: its features
are pretrained to localize objects across scales, which is close to what
stereo matching needs from an encoder, and the encoder ablation of
Chapter 4 confirmed that this backbone beats both a hand-built GhostConv
stack and its own smaller sibling on every accuracy metric at a
two-millisecond cost. The detection neck and head above the one-sixteenth
stage are removed at construction rather than merely unused, because
dead blocks would waste parameters and forward compute. The two views
are concatenated along the batch dimension and pass through a single
shared invocation, which guarantees identical weights for both images.
A separate context encoder
processes the left image alone. It applies a strided convolution followed
by two GhostConv blocks with squeeze-and-excitation, and projects to a
32-channel quarter-resolution context map that is pooled to the coarser
scales. Matching evidence and recurrent context are kept in separate
streams for a specific reason. When the matching features also serve as
context, the recurrent unit cannot separate appearance from
correspondence; the RAFT-Stereo lineage \cite{lipson2021raftstereo}
established that a dedicated context path improves refinement, and the
same separation is adopted here.

\begin{figure}[!htbp]
\centering
\includegraphics[width=0.95\textwidth]{fig_3_2_encoders_preview}
\caption{\textbf{Feature and context encoders.} The shared YOLO26s encoder processes both views and emits matching features at three scales (f4 128, f8 256, f16 256 channels); a separate GhostConv context encoder sees only the left image and produces a 32-channel quarter-resolution context map, pooled to the coarser scales.}
\label{fig:encoders}
\end{figure}

\section{Tile Initialization and Cost Volume}
% Source: FINAL_MODEL_ARCHITECTURE section 6. Introduce the tile state
% T = (d, sx, sy, h, c) here.
To seed the coarse-to-fine process, the tile initialization module builds
a group-wise correlation volume at one sixteenth resolution. For each of
eight feature groups and each integer disparity, the matching score is the
mean channel-wise product between the left features and the horizontally
shifted right features, given by \eqref{eq:gwc}. A small three-dimensional
convolutional network with an eight, sixteen, sixteen, one channel
schedule regularizes the volume, and a softmax along the disparity axis
yields a probability distribution. The initial disparity and confidence
follow from the soft-argmax and the distribution peak in
\eqref{eq:softargmax_init}. The plane slopes are initialized to zero and
the hidden feature is replaced by the coarsest context map. The 24
hypotheses at this scale span roughly 0 to 368 full-resolution pixels.

\begin{equation}
C(g,d,y,x) \;=\; \frac{1}{C/G}\sum_{c=1}^{C/G}
  f_{L}^{\,g,c}(y,x)\; f_{R}^{\,g,c}(y,\,x-d)
\label{eq:gwc}
\end{equation}
where $G = 8$ is the number of feature groups, $C$ the channel count,
$d \in \{0, \dots, 23\}$ the integer disparity hypotheses at one sixteenth
resolution, and $f_{L}^{\,g,c}$, $f_{R}^{\,g,c}$ the $c$-th channel of
group $g$ of the left and right feature maps.

\begin{equation}
d_{0} \;=\; \sum_{d} p(d)\, d, \qquad
c_{0} \;=\; \max_{d}\, p(d)
\label{eq:softargmax_init}
\end{equation}
where $p(d)$ is the softmax of the regularized cost volume along the
disparity dimension, $d_0$ the expected disparity, and $c_0$ the peak
probability that serves as an initial confidence.

\begin{figure}[!htbp]
\centering
\includegraphics[width=0.95\textwidth]{fig_3_3_tile_init_preview}
\caption{\textbf{Tile-state initialization at 1/16 resolution.} A group-wise correlation volume (8 groups, 24 hypotheses) is regularized by a small 3-D CNN; soft-argmax and the distribution maximum yield the initial disparity and confidence, and the context map replaces the zero hidden state.}
\label{fig:tile_init}
\end{figure}

\section{Recurrent Tile Refinement}
% Source: FINAL_MODEL_ARCHITECTURE section 7. Separate unshared
% TileRefine modules at 1/16, 1/8, 1/4 (channel counts differ); warp +
% local correlation lookup (half range 2, five offsets); the gated
% residual head (four sigmoid gates); softplus non-negativity; schedule
% 2/3/3 = eight updates.
The tile state is refined by separate refinement modules at one
sixteenth, one eighth, and one quarter resolution; their weights are not
shared because the matching-feature channel counts differ across scales.
At each iteration the right feature map is warped toward the left view by
the current disparity, and a differentiable local correlation is
evaluated over five offsets within a half-range of two pixels. The
recurrent input assembles the left features, the warped right features,
the current state, the local correlation, and the context, as in
\eqref{eq:gru_input}. A convolutional gated recurrent unit updates the
persistent hidden state through the gate equations of \eqref{eq:convgru}.
A two-layer head with 48 hidden channels predicts residuals for
disparity, both plane slopes, and confidence, and a second learned gate
emits four independent sigmoid weights that scale each residual.
Disparity is constrained to remain non-negative through a softplus. The
schedule runs two iterations at one sixteenth, three at one eighth, and
three at one quarter, for eight learned updates in total.

\begin{equation}
x \;=\; \bigl[\, f_{L},\; \mathcal{W}(f_{R};\,d),\; d,\; s_{x},\; s_{y},\;
  c,\; \mathrm{corr},\; \mathrm{ctx} \,\bigr]
\label{eq:gru_input}
\end{equation}
where $\mathcal{W}(f_{R}; d)$ denotes the right features warped by the
current disparity, $\mathrm{corr}$ the five-offset local correlation, and
$\mathrm{ctx}$ the scale-matched context map.

\begin{equation}
\begin{aligned}
z &= \sigma\!\bigl(\mathrm{Conv}([h,\,x])\bigr), &
r &= \sigma\!\bigl(\mathrm{Conv}([h,\,x])\bigr),\\
q &= \tanh\!\bigl(\mathrm{Conv}([r \odot h,\,x])\bigr), &
h' &= (1-z) \odot h + z \odot q,
\end{aligned}
\label{eq:convgru}
\end{equation}
where $h$ is the 32-channel persistent hidden state, $z$ and $r$ the
update and reset gates, $q$ the candidate state, and $\odot$ the
element-wise product.

\begin{figure}[!htbp]
\centering
\includegraphics[width=\textwidth]{fig_3_4_refinement_preview}
\caption{\textbf{One recurrent refinement iteration.} The right features are warped by the current disparity, a five-offset local correlation is evaluated, and the assembled input updates the persistent hidden state through a convolutional GRU; four sigmoid gates control the residual updates of disparity, both slopes, and confidence.}
\label{fig:refinement}
\end{figure}

\section{Plane-Aware Cross-Scale Propagation}
% Source: FINAL_MODEL_ARCHITECTURE section 8. Slopes are USED, not just
% stored; contrast with plain bilinear upsampling.
Between adjacent scales, the propagation module resamples the disparity,
slopes, hidden feature, and confidence to the finer grid. Disparity is
multiplied by two because its pixel unit changes with resolution. Unlike
plain bilinear interpolation, the propagated disparity also uses the local
plane slopes to place the four child pixels around each parent tile,
following \eqref{eq:plane_prop}. This is the tile-plane component of the
design, in the manner of HITNet \cite{tankovich2021hitnet}; it preserves
a local slanted-surface hypothesis while moving to a finer grid, which
keeps sub-pixel structure without rebuilding a cost volume at every level.

\begin{equation}
d_{\mathrm{child}} \;=\; 2\,\mathcal{B}(d_{\mathrm{parent}})
  \;+\; 2\,s_{x}\,\Delta x \;+\; 2\,s_{y}\,\Delta y,
\qquad \Delta x, \Delta y \in \{-0.25,\, +0.25\}
\label{eq:plane_prop}
\end{equation}
where $\mathcal{B}$ is bilinear resampling of the parent disparity, $s_x$
and $s_y$ the local plane slopes, and $\Delta x, \Delta y \in \{-0.25,
+0.25\}$ the child offsets; the factor of two rescales the disparity to
the finer grid.

\begin{figure}[!htbp]
\centering
\includegraphics[width=\textwidth]{fig_3_5_plane_prop_preview}
\caption{\textbf{Plane-aware propagation between scales.} Each parent tile spawns four children whose disparities follow the local plane defined by the slopes, per \eqref{eq:plane_prop}, instead of plain bilinear interpolation.}
\label{fig:plane_prop}
\end{figure}

\section{Geometry Encoding Volume and Fail-Soft Fusion}
% Source: FINAL_MODEL_ARCHITECTURE section 9 + the narrow-band variant
% (the trained configuration): the volume samples 33 hypotheses in a
% +-16 band around the current tile disparity (efficiency finding F3,
% validated accuracy-neutral). Three 3x3x3 convolutions, 16 channels.
The geometry encoding volume adds an independent global proposal at one
quarter resolution, in the spirit of IGEV-Stereo \cite{xu2023igev}. In the
trained configuration it is narrow: rather than the full 64 disparity
planes of the base design, it samples 33 hypotheses in a band of plus or
minus 16 around the current tile disparity, a change validated in
Chapter 4 as accuracy-neutral while removing roughly a quarter of the
model's multiply-accumulate operations. The band is not as restrictive
as it sounds: 16 quarter-resolution planes correspond to 64
full-resolution pixels on either side of the current estimate, and since
the coarse initialization has already localized the disparity, the
band's job is correction rather than search. The full-range volume
spends most of its computation scoring hypotheses the tile state has
already ruled out, which is why removing them costs nothing. The volume is regularized by
three three-dimensional convolutions with 16 hidden channels, and the
softmax distribution yields the expected disparity, a confidence, and a
projected geometry feature, as in \eqref{eq:gev_expect}. The proposal is
not substituted unconditionally. A learned gate takes the context, the
geometry feature, both confidences, the absolute disagreement, and the
current disparity, and produces a scalar weight through
\eqref{eq:failsoft}. The gate's final bias is initialized to minus four,
so the initial weight is approximately 0.018 and training begins close to
the stable recurrent baseline; the network then learns where the geometry
proposal helps. This fail-soft design keeps the global branch from
corrupting regions the recurrent path already resolves well.

\begin{equation}
d_{\mathrm{gev}} = \sum_{d} p(d)\, d, \qquad
c_{\mathrm{gev}} = \max_{d}\, p(d), \qquad
g_{\mathrm{gev}} = \phi\Bigl(\sum_{d} G(d)\, p(d)\Bigr)
\label{eq:gev_expect}
\end{equation}
where $p(d)$ is the softmax over the regularized quarter-resolution
volume, $G(d)$ the regularized three-dimensional geometry feature, and
$\phi$ a projection to a 16-channel two-dimensional expected-geometry map.

\begin{equation}
\begin{aligned}
w &= \sigma\!\Bigl(F\bigl([\mathrm{ctx}_{4},\, g_{\mathrm{gev}},\, c,\,
  c_{\mathrm{gev}},\, \lvert d_{\mathrm{gev}} - d \rvert,\, d]\bigr)\Bigr),\\
d_{\mathrm{fused}} &= \mathrm{softplus}\bigl(d + w\,(d_{\mathrm{gev}} - d)\bigr),
\end{aligned}
\label{eq:failsoft}
\end{equation}
where $F$ is the fusion network whose final-layer bias is initialized to
$-4$, $c$ and $c_{\mathrm{gev}}$ the tile and geometry confidences, and
$w$ the blend weight; the slopes are attenuated by $(1 - w)$ and the
confidence becomes the maximum of the two estimates.

\begin{figure}[!htbp]
\centering
\includegraphics[width=\textwidth]{fig_3_6_gev_fusion_preview}
\caption{\textbf{Narrow-band geometry encoding volume and fail-soft fusion.} The volume samples 33 hypotheses in a $\pm 16$ band around the current tile disparity; a learned gate with bias initialized to $-4$ blends the GEV proposal into the tile estimate only where it helps.}
\label{fig:gev_fusion}
\end{figure}

\section{Plane Rendering and Learned Convex Upsampling}
% Source: FINAL_MODEL_ARCHITECTURE section 10 + the plane-render fix
% (blur-fix ablation 20260703_blurfix_n500: bad-0.5 -22.5 percent,
% selected over bimodal and bundle1; composition was WORSE than plane
% alone: report that in Ch4, cite here).
The final one-quarter disparity is upsampled to full resolution in two
learned two-times stages. The first stage predicts masks from the
one-quarter left features and the second from the one-half left features.
For every output subpixel, a mask predicts nine weights over the
three-by-three coarse neighborhood, and a softmax makes the weights
non-negative and sum to one, so the upsampled disparity is a convex
combination of nearby coarse values, following the convex upsampling of
RAFT-Stereo \cite{lipson2021raftstereo}. Before upsampling, a
plane-rendering step reconstructs each tile from its disparity and slopes
and applies an edge gate; this was the decisive fix in the
boundary-sharpness study of Chapter 4, where it reduced the bad-0.5 rate
by roughly 22 percent relative to the control. Convex upsampling can
select different source neighbors near a boundary, but it cannot invent a
disparity hypothesis absent from the one-quarter prediction, so thin and
distant structures remain a known limitation.

\begin{figure}[!htbp]
\centering
\includegraphics[width=0.92\textwidth]{fig_3_7_upsample_preview}
\caption{\textbf{Learned convex upsampling.} Every output subpixel is a softmax-weighted convex combination of its $3\times 3$ coarse neighborhood, with masks predicted from the left features; two 2$\times$ stages recover full resolution.}
\label{fig:upsample}
\end{figure}

\section{Training Objective}
% Source: FINAL_MODEL_ARCHITECTURE section 12 + trainer loss_fn.
% Valid pixels: finite GT in (0, 320) px. Lower scales resized to full
% resolution and scaled before comparison.
The network is trained with the multi-scale supervised objective of
\eqref{eq:loss}. A pixel counts as valid when its ground-truth disparity
is finite and lies between 0 and 320 pixels. Each lower-scale prediction
is resized to full resolution and multiplied by its scale factor before
the comparison. The objective weights an L1 term at every scale most
heavily at full resolution, adds a matching L1 term on the
geometry-volume output, and augments the full-resolution prediction with
a Sobel gradient-consistency term, a threshold hinge stack, a KITTI-style
D1-relative hinge, an edge-aware smoothness term, and a gated
slant-supervision term that trains the plane slopes only where they are
trustworthy. The slant term pairs with plane rendering: it makes the
slopes accurate enough for the rendering step to sharpen boundaries. The
composition of the objective is itself an experimental result rather
than a guess. The loss ablation of Chapter 4 tested nine formulations
and found that a plain L1 under-serves the outlier metrics, that hinge
terms on a single threshold overfit that threshold, and that the
threshold stack combined with the D1 hinge was the only formulation to
win the end-point error and the outlier rates together; the production
objective descends from that winner, with the gradient term retained for
its sub-pixel benefit and the smoothness term kept at a deliberately
small weight so it regularizes without blurring.

\begin{equation}
\begin{aligned}
\mathcal{L} \;=\;& 1.00\,\mathcal{L}_{1}(d_{\mathrm{full}})
 + 0.50\,\mathcal{L}_{1}(d_{1/2})
 + 0.30\,\mathcal{L}_{1}(d_{1/4})
 + 0.20\,\mathcal{L}_{1}(d_{1/8})
 + 0.10\,\mathcal{L}_{1}(d_{1/16})\\
&+ 0.15\,\mathcal{L}_{1}(d_{\mathrm{gev}})
 + 0.50\,\mathcal{L}_{\mathrm{grad}}
 + 0.20\,\mathcal{L}_{\mathrm{thr}}
 + 0.20\,\mathcal{L}_{\mathrm{D1}}
 + 0.02\,\mathcal{L}_{\mathrm{smooth}}
 + 0.30\,\mathcal{L}_{\mathrm{slant}}
\end{aligned}
\label{eq:loss}
\end{equation}
where $\mathcal{L}_1$ is the masked mean absolute disparity error,
$\mathcal{L}_{\mathrm{grad}}$ a Sobel gradient-consistency term,
$\mathcal{L}_{\mathrm{thr}}$ a hinge over the error thresholds
$\{0.5, 1, 2, 3\}$ pixels, $\mathcal{L}_{\mathrm{D1}}$ the D1-relative
hinge, $\mathcal{L}_{\mathrm{smooth}}$ an edge-aware smoothness term, and
$\mathcal{L}_{\mathrm{slant}}$ the gated slant supervision.

\begin{figure}[!htbp]
\centering
\includegraphics[width=\textwidth]{fig_3_8_supervision_preview}
\caption{\textbf{Multi-scale supervision.} Six disparity outputs receive weighted L1 terms against the ground truth per \eqref{eq:loss}; the full-resolution prediction additionally receives the gradient, threshold, D1, smoothness, and gated slant terms.}
\label{fig:supervision}
\end{figure}

\section{Datasets and Training Protocol}
% Sources: run meta.json args block; 20260704_native_vs_resize_n500
% comparison.md; the grand-comparison augmentation evidence.
% Content: SceneFlow canonical split (35,454 train = FT3D-TRAIN + Monkaa
% + Driving finalpass; 4,370 test = FT3D-TEST); native 384x640 co-located
% random crops (cite the native-vs-resize ablation, Ch4); OpenStereo
% augmentation triplet (asymmetric color jitter, right-only random erase,
% random scale); OneCycle 60k steps peak 8e-4, batch 32, bf16, A100;
% input contract [0,1]; validation on a fixed 400-pair FT3D-TEST subset
% at the native 960x540 pad-16 axis.
The model is trained on the full SceneFlow finalpass dataset. The
canonical split contains 35,454 training pairs, drawn from the
FlyingThings3D training set together with Monkaa and Driving, and the
FlyingThings3D test set of 4,370 pairs is held out for evaluation; only
FlyingThings3D provides an official test split, so the test set is drawn
from it alone. Training samples random 384 by 640 windows that are
co-located in both views from the native 960 by 540 frames, which
preserves the native disparity magnitudes. Chapter 4 shows that this
native-crop protocol substantially outperforms the legacy protocol of
downscaling the whole frame. Three augmentations from the OpenStereo
recipe are applied: asymmetric color jitter, right-image-only random
erasing, and random scaling. Optimization uses AdamW with a OneCycle
schedule over 60,000 steps at a peak learning rate of $8 \times 10^{-4}$,
a batch of 32, and bfloat16 mixed precision, on a single A100 GPU. The
input is normalized to the unit interval at the model boundary.
Validation uses a fixed 400-pair subset of the test set at the native
960 by 540 axis, evaluated every thousand steps, with a resized-axis
measurement logged alongside. Table \ref{tab:train_config} lists the
configuration.

\begin{figure}[!htbp]
\centering
\includegraphics[width=\textwidth]{fig_3_10_input_protocol}
\caption{\textbf{Native-crop input protocol.} Training samples random $384\times 640$ co-located windows of the native $960\times 540$ frame (left), preserving native disparity magnitudes; the legacy protocol (right) downscales the whole frame and shrinks every disparity by the same factor.}
\label{fig:input_protocol}
\end{figure}

\begin{table}[!htbp]
\centering
\caption{\textbf{Training configuration of the production run.} Values verbatim from the run's configuration record (meta.json of run 20260704).}
\label{tab:train_config}
\small
\begin{tabular}{ll}
\toprule
Element & Value\\
\midrule
Optimizer & AdamW, weight decay $10^{-5}$\\
Schedule & OneCycle, 60{,}000 steps, peak LR $8\times 10^{-4}$\\
Batch size / precision & 32 / bfloat16 mixed\\
Input protocol & random native $384\times 640$ co-located crops\\
Augmentation & asymmetric color jitter, right-only erase, random scale\\
Training data & full SceneFlow finalpass, 35{,}454 pairs\\
Validation & fixed 400-pair FlyingThings3D subset, every 1k steps\\
Hardware & one NVIDIA A100 80 GB (cloud)\\
\bottomrule
\end{tabular}
\end{table}

\section{Design-Space Exploration Methodology}
% Sources: model/benchmarks/OVERFIT_METHODOLOGY.md; the ablation-study-
% expert protocol. Content: matched-overfit harness (fixed pairs, fixed
% seed, full 8-metric reporting), controlled single-knob ablations,
% pre-registered win criteria, cross-dataset gate after every training
% run. This section is methodology (PO4 setup); its RESULTS are Ch4.
The architectural choices in this chapter were made under a controlled
matched-overfit methodology rather than by intuition. Each candidate is
trained on a fixed set of SceneFlow pairs with a fixed seed, deterministic
convolution kernels, an identical step budget, and an identical evaluation
protocol, and all eight stereo metrics together with an inference-latency
benchmark are reported for every arm. Single knobs are varied one at a
time, the winning criterion is registered before a run, and every
candidate is checked against a held-out cross-dataset evaluation, so that
an in-domain gain that fails to transfer is not mistaken for progress. The
results of these studies appear in Chapter 4; this section states the
protocol that produced them.

\section{Optimization of the Network for Edge Inference}
% Source: THESIS_PLAN section 3 (F1-F8 findings table). This is thesis
% content evidencing PO3/PO5: cost-volume zero-copy views, grid-sample
% batching, static-channel hoisting, head fusion, dead-code removal,
% narrow GEV band; equivalence-proof discipline (bitwise or matched A/B
% per fix); export blockers identified (F.unfold, GroupNorm, Conv3d) and
% their mitigations.
Reaching the latency target required an optimization pass over the
network's inference graph, treated as an engineering task in its own
right. Six changes were applied. The cost-volume construction loops were
rewritten to use padded zero-copy views instead of allocating and zeroing
a feature map for every disparity bin. The grid-sample calls were reduced
from 48 to eight by batching the lookups and caching the sampling grids.
Static channels of the recurrent input, which are re-convolved every
iteration, were hoisted once per scale, following the precomputed-context
trick of RAFT-Stereo \cite{lipson2021raftstereo}. The gate convolutions
and the four output heads, which share inputs, were fused. Dead
inference-path code was removed. Finally, the geometry volume was narrowed
to a band around the tile disparity. Each change was proven equivalent
before adoption, either bitwise or through a matched A/B comparison, and
together they cut the full-precision latency by a factor of 1.74. The pass
also surfaced three export blockers for eight-bit deployment, namely the
unfold operation in convex upsampling, group normalization on some
accelerator toolchains, and three-dimensional convolution on neural
processing units, each of which has a documented replacement.
Table \ref{tab:efficiency_findings} summarizes the changes.

\begin{table}[!htbp]
\centering
\caption{\textbf{Inference optimizations and their equivalence status.} Savings measured at $384\times 640$ on RTX 3050; equivalence class states how each change was proven safe before adoption.}
\label{tab:efficiency_findings}
\small
\begin{tabular}{p{5.6cm}p{4.2cm}l}
\toprule
Optimization & Saving & Equivalence\\
\midrule
Zero-copy padded views in cost-volume loops & 5 to 15\% latency & bitwise\\
Batched correlation lookups, cached grids & 3 to 8\% latency & identical math\\
Static GRU input channels hoisted per scale & 8 to 10\% of MACs & exact\\
Fused gate convolutions and output heads & 1 to 3\% latency & bitwise\\
Dead inference-path code removed & memory + launches & exact\\
Narrow GEV band ($\pm 16$ around tile $d$) & 25 to 30\% of MACs & matched A/B (Ch.~4)\\
\midrule
Combined & fp32 61.4 ms vs 106.7 ms base ($1.74\times$) & \\
\bottomrule
\end{tabular}
\end{table}

\begin{figure}[!htbp]
\centering
\includegraphics[width=0.8\textwidth]{fig_3_11_param_budget}
\caption{\textbf{Parameter budget of the trained model.} Module split computed directly from the trained checkpoint's state dictionary; the shared encoder and the three refinement stages hold 95 percent of the 2.96 M parameters.}
\label{fig:param_budget}
\end{figure}

\section{Evaluation Metrics}
% Source: FINAL_MODEL_ARCHITECTURE section 16.
Accuracy is reported with the standard stereo metrics defined in
\eqref{eq:metrics}. The end-point error is the mean absolute disparity
error over valid pixels; the bad-$t$ rate is the percentage of valid
pixels whose absolute error exceeds $t$ pixels; and the D1-all rate is the
KITTI outlier definition. The root-mean-square error and the median
absolute error are reported alongside, the former to expose large outliers
and the latter as an outlier-robust summary. Each image contributes its
own averaged metric, and the scores are averaged over images rather than
by pooling every pixel globally. Latency is measured at batch one and
384 by 640 resolution with ten warm-up passes and a hundred timed passes,
and memory as the peak inference working set; the device is stated for
every reported measurement.

\begin{equation}
\begin{aligned}
\mathrm{EPE} &= \frac{1}{\lvert V \rvert} \sum_{p \in V}
  \bigl\lvert \hat{d}(p) - d^{*}(p) \bigr\rvert, \qquad
\mathrm{bad}\text{-}t = \frac{100}{\lvert V \rvert}\,
  \Bigl\lvert \bigl\{ p \in V :
  \lvert \hat{d}(p) - d^{*}(p) \rvert > t \bigr\} \Bigr\rvert,\\
\mathrm{D1} &= \frac{100}{\lvert V \rvert}\, \Bigl\lvert \bigl\{ p \in V :
  \lvert \hat{d}(p) - d^{*}(p) \rvert > 3 \;\wedge\;
  \lvert \hat{d}(p) - d^{*}(p) \rvert > 0.05\, d^{*}(p) \bigr\} \Bigr\rvert
\end{aligned}
\label{eq:metrics}
\end{equation}
where $V$ is the set of valid pixels, $\hat{d}$ the predicted disparity,
$d^{*}$ the ground truth, and $t \in \{0.5, 1, 2, 3\}$ pixels the bad-rate
thresholds.

\section{Summary of Design Decisions}
Every load-bearing decision in this chapter is tied to a controlled
study in Chapter 4 or to published evidence in Chapter 2. Table
\ref{tab:design_evidence} collects the mapping so that the design can be
audited decision by decision.

\begin{table}[!htbp]
\centering
\caption{\textbf{Design decisions and their evidence.} Each architectural or training choice with the study or literature that justifies it.}
\label{tab:design_evidence}
\small
\begin{tabular}{p{4.6cm}p{7.6cm}}
\toprule
Decision & Evidence\\
\midrule
YOLO26s shared encoder & encoder ablation, Table \ref{tab:ablation_encoder}: best EPE of three encoders at +2 ms\\
Separate left context encoder & RAFT-Stereo lineage \cite{lipson2021raftstereo}; context-gated variants win the family ablation, Table \ref{tab:ablation_gev4line}\\
Group-wise cost at one coarse scale & GwcNet \cite{guo2019gwcnet} graceful degradation under small aggregation budgets\\
Iterated coarse-to-fine refinement & refinement ablation, Table \ref{tab:ablation_refinement}: single-pass collapses, finer scales add cost without gain\\
Local $\pm 2$ correlation lookup & CREStereo \cite{li2022crestereo}: local search beats all-pairs; also confers rectification tolerance, Section \ref{sec:rectification}\\
Slanted tile planes & HITNet \cite{tankovich2021hitnet}; plane rendering wins the blur study, Table \ref{tab:ablation_blurfix}\\
Narrow GEV band + fail-soft gate & family ablation best EPE, Table \ref{tab:ablation_gev4line}; narrowing accuracy-neutral, Table \ref{tab:ablation_efficiency}\\
Rejected: selective recurrent unit & family ablation, Table \ref{tab:ablation_gev4line}: EPE worsens on the compact chassis\\
Rejected: fresh 1/4 initialization & family ablation, Table \ref{tab:ablation_gev4line}: EPE and bad-1 worsen\\
Rejected: fix compositions & blur study, Table \ref{tab:ablation_blurfix}: worse than plane rendering alone\\
Threshold-stack + D1 objective & loss ablation, Table \ref{tab:ablation_loss}: wins EPE, bad-2, D1 together\\
Native-crop input protocol & input-protocol study, Table \ref{tab:ablation_native}: native axis EPE 6.67 to 2.87 at no resized-axis cost\\
Augmentation triplet & matched comparison: best val EPE 2.778 to 1.921\\
OneCycle 60k schedule & training dynamics, Chapter 4: validation floor reached before the budget ends\\
\bottomrule
\end{tabular}
\end{table}

\section{From Disparity to 3D Reconstruction}
\label{sec:reconstruction_method}
The disparity map is not the end product of the pipeline; it is the
carrier of the three-dimensional structure of the scene. Once the
disparity is known, every valid pixel back-projects to a metric
three-dimensional point through the pinhole relations
\begin{equation}
Z = \frac{f\,B}{d}, \qquad
X = \frac{(u - c_x)\,Z}{f}, \qquad
Y = \frac{(v - c_y)\,Z}{f}
\label{eq:backprojection}
\end{equation}
where $(u, v)$ are the pixel coordinates, $(c_x, c_y)$ the principal
point, $f$ the focal length in pixels, $B$ the stereo baseline, and
$(X, Y, Z)$ the recovered point in the camera frame. Coloring each point
with the corresponding left-image pixel yields a dense colored point
cloud of the scene. For the reconstructions reported in Chapter 4, the
cloud is post-processed with two standard operations: a statistical
outlier filter removes the isolated flying points that appear along depth
discontinuities, and a voxel downsampling makes the point density
uniform. The stereo rig's calibrated values, a focal length of
approximately 1005 pixels at the native 1280-pixel width and a baseline
of 52 millimeters, scale to the inference resolution by the width ratio.
This reconstruction path is what the thesis title refers to: the network
provides the dense disparity, and the calibrated geometry turns it into a
metric three-dimensional model of the scene.

\section{Implementation and Deployment Environment}
% Sources: meta.json; THESIS_PLAN section 2A. PyTorch 2.11 cu128; Modal
% A100-80GB training, T4 evaluation; RTX 3050 local bench; CCB stereo
% camera rig; target device Jetson Orin Nano (measured on-device 36.3 ms).
The model is implemented in PyTorch 2.11 with CUDA 12.8. Training runs on
a single cloud A100 with 80 gigabytes of memory, ablations and evaluations
on a cloud T4, and the local latency benchmark on an RTX 3050 laptop GPU
with roughly 4 gigabytes of memory. Real-camera pairs come from a
calibrated binocular rig that streams a 2560 by 720 side-by-side image.
The intended deployment target is the Jetson Orin Nano, on which the
optimized model was measured directly; the on-device latency is reported
in Chapter 4.
